\documentclass[12pt]{article}
\usepackage{amsmath}
\usepackage{amssymb}
\usepackage{geometry}
\geometry{a4paper, margin=0.5in}

\begin{document}

\section*{Part A: LU Factorization}
\hrulefill
\vspace{1em}

% BEGIN QUESTION 1
\subsection*{Question 1}
Find the LU factorization of the matrix:
$$
A = \begin{bmatrix}
3 & 2 & 1 \\
6 & 5 & 4 \\
9 & 8 & 7
\end{bmatrix}
$$

\subsection*{Solution}

We perform Gaussian elimination to find $L$ and $U$.

Starting with:
\[
	A = \begin{bmatrix}
		3 & 2 & 1 \\
		6 & 5 & 4 \\
		9 & 8 & 7
	\end{bmatrix}
\]

$R_2 - 2R_1$:
\[
	\begin{bmatrix}
		3 & 2 & 1 \\
		0 & 1 & 2 \\
		9 & 8 & 7
	\end{bmatrix}, \quad l_{21} = 2
\]

$R_3 - 3R_1$:
\[
	\begin{bmatrix}
		3 & 2 & 1 \\
		0 & 1 & 2 \\
		0 & 2 & 4
	\end{bmatrix}, \quad l_{31} = 3
\]

$R_3 - 2R_2$:
\[
	U = \begin{bmatrix}
		3 & 2 & 1 \\
		0 & 1 & 2 \\
		0 & 0 & 0
	\end{bmatrix}, \quad l_{32} = 2
\]

Using the multipliers:
\[
	L = \begin{bmatrix}
		1 & 0 & 0 \\
		2 & 1 & 0 \\
		3 & 2 & 1
	\end{bmatrix}
\]

Therefore, the LU factorization is:
\[
	{A = LU = \begin{bmatrix}
		1 & 0 & 0 \\
		2 & 1 & 0 \\
		3 & 2 & 1
	\end{bmatrix} \begin{bmatrix}
		3 & 2 & 1 \\
		0 & 1 & 2 \\
		0 & 0 & 0
	\end{bmatrix}}
\]

Verification:
\begin{align*}
	LU &= \begin{bmatrix}
		1 & 0 & 0 \\
		2 & 1 & 0 \\
		3 & 2 & 1
	\end{bmatrix} \begin{bmatrix}
		3 & 2 & 1 \\
		0 & 1 & 2 \\
		0 & 0 & 0
	\end{bmatrix} \\
	&= \begin{bmatrix}
		3 & 2 & 1 \\
		6 & 5 & 4 \\
		9 & 8 & 7
	\end{bmatrix} = A
\end{align*}
% END QUESTION 1

% BEGIN QUESTION 2
\subsection*{Question 2}
Given the matrix:
$$
B = \begin{bmatrix}
4 & 1 & 2 \\
8 & 5 & 3 \\
12 & 7 & 11
\end{bmatrix}
$$
Find its LU decomposition and then solve the system $Bx = \begin{bmatrix} 5 \\ 13 \\ 25 \end{bmatrix}$.

\subsection*{Solution}

Starting with:
\[
	B = \begin{bmatrix}
		4 & 1 & 2 \\
		8 & 5 & 3 \\
		12 & 7 & 11
	\end{bmatrix}
\]

$R_2 - 2R_1$:
\[
	\begin{bmatrix}
		4 & 1 & 2 \\
		0 & 3 & -1 \\
		12 & 7 & 11
	\end{bmatrix}, \quad l_{21} = 2
\]

$R_3 - 3R_1$:
\[
	\begin{bmatrix}
		4 & 1 & 2 \\
		0 & 3 & -1 \\
		0 & 4 & 5
	\end{bmatrix}, \quad l_{31} = 3
\]

$R_3 - \frac{4}{3}R_2$:
\[
	U = \begin{bmatrix}
		4 & 1 & 2 \\
		0 & 3 & -1 \\
		0 & 0 & \frac{19}{3}
	\end{bmatrix}, \quad l_{32} = \frac{4}{3}
\]

\[
	L = \begin{bmatrix}
		1 & 0 & 0 \\
		2 & 1 & 0 \\
		3 & \frac{4}{3} & 1
	\end{bmatrix}
\]

Therefore, the LU factorization is:
\[
	B = LU = \begin{bmatrix}
		1 & 0 & 0 \\
		2 & 1 & 0 \\
		3 & \frac{4}{3} & 1
	\end{bmatrix} \begin{bmatrix}
		4 & 1 & 2 \\
		0 & 3 & -1 \\
		0 & 0 & \frac{19}{3}
	\end{bmatrix}
\]

Solve $Bx = \begin{bmatrix} 5 \\ 13 \\ 25 \end{bmatrix}$ using $Ly = b$:
\[
	\begin{bmatrix}
		1 & 0 & 0 \\
		2 & 1 & 0 \\
		3 & \frac{4}{3} & 1
	\end{bmatrix} \begin{bmatrix}
		y_1 \\ y_2 \\ y_3
	\end{bmatrix} = \begin{bmatrix}
		5 \\ 13 \\ 25
	\end{bmatrix}
\]

\begin{align*}
	y_1 &= 5 \\
	2y_1 + y_2 &= 13 \Rightarrow y_2 = 3 \\
	3y_1 + \frac{4}{3}y_2 + y_3 &= 25 \Rightarrow y_3 = 6
\end{align*}

Solve $Ux = y$:
\[
	\begin{bmatrix}
		4 & 1 & 2 \\
		0 & 3 & -1 \\
		0 & 0 & \frac{19}{3}
	\end{bmatrix} \begin{bmatrix}
		x_1 \\ x_2 \\ x_3
	\end{bmatrix} = \begin{bmatrix}
		5 \\ 3 \\ 6
	\end{bmatrix}
\]

\begin{align*}
	\frac{19}{3}x_3 &= 6 \Rightarrow x_3 = \frac{18}{19} \\
	3x_2 - x_3 &= 3 \Rightarrow x_2 = \frac{25}{19} \\
	4x_1 + x_2 + 2x_3 &= 5 \Rightarrow x_1 = \frac{17}{38}
\end{align*}

Therefore: $x_1 = \frac{17}{38}, x_2 = \frac{25}{19}, x_3 = \frac{18}{19}$
% END QUESTION 2

% BEGIN QUESTION 3
\subsection*{Question 3}
Determine if the following matrix has an LU factorization. If yes, find it:
$$
C = \begin{bmatrix}
1 & 2 & 3 \\
2 & 4 & 7 \\
3 & 6 & 10
\end{bmatrix}
$$

\subsection*{Solution}

Starting with:
\[
	C = \begin{bmatrix}
		1 & 2 & 3 \\
		2 & 4 & 7 \\
		3 & 6 & 10
	\end{bmatrix}
\]

$R_2 \leftarrow R_2 - 2R_1$:
\[
	\begin{bmatrix}
		1 & 2 & 3 \\
		0 & 0 & 1 \\
		3 & 6 & 10
	\end{bmatrix}, \quad l_{21} = 2
\]

Zero pivot at position $(2,2)$. The matrix does not have a standard LU factorization.
% END QUESTION 3

\vspace{2em}
\section*{Part B: Linear Transformations}
\hrulefill
\vspace{1em}

% BEGIN QUESTION 4
\subsection*{Question 4}
Consider the linear transformation $T : \mathbb{R}^{2} \rightarrow \mathbb{R}^{2}$ defined by:
$$
T(x,y) = (2x + 3y, x - 4y)
$$
Find the standard matrix of T and determine if it is invertible.

\subsection*{Solution}

Apply $T$ to standard basis vectors:
\begin{align*}
	T(1,0) &= (2, 1) \\
	T(0,1) &= (3, -4)
\end{align*}

Standard matrix:
\[
	A = \begin{bmatrix}
		2 & 3 \\
		1 & -4
	\end{bmatrix}
\]

Calculate determinant:
\[
	\det(A) = \begin{vmatrix}
		2 & 3 \\
		1 & -4
	\end{vmatrix} = 2(-4) - 3(1) = -8 - 3 = -11
\]

Since $\det(A) = -11 \neq 0$, the transformation is invertible.

Inverse matrix:
\[
	A^{-1} = \frac{1}{-11}\begin{bmatrix}
		-4 & -3 \\
		-1 & 2
	\end{bmatrix} = \begin{bmatrix}
		\frac{4}{11} & \frac{3}{11} \\
		\frac{1}{11} & -\frac{2}{11}
	\end{bmatrix}
\]
% END QUESTION 4

% BEGIN QUESTION 5
\subsection*{Question 5}
Let $T : \mathbb{R}^{3} \rightarrow \mathbb{R}^{2}$ be a linear transformation such that:
\begin{align*}
T(1,0,0) &= (2,1) \\
T(0,1,0) &= (-1,3) \\
T(0,0,1) &= (4,-2)
\end{align*}
Find $T(3,-2,5)$ and determine the kernel of T.

\subsection*{Solution}

Standard matrix from given basis images:
\[
	A = \begin{bmatrix}
		2 & -1 & 4 \\
		1 & 3 & -2
	\end{bmatrix}
\]

Find $T(3,-2,5)$:
\begin{align*}
	T(3,-2,5) &= \begin{bmatrix}
		2 & -1 & 4 \\
		1 & 3 & -2
	\end{bmatrix} \begin{bmatrix}
		3 \\ -2 \\ 5
	\end{bmatrix} \\
	&= \begin{bmatrix}
		6 + 2 + 20 \\
		3 - 6 - 10
	\end{bmatrix} \\
	&= \begin{bmatrix}
		28 \\ -13
	\end{bmatrix}
\end{align*}

Therefore: $T(3,-2,5) = (28, -13)$

Find kernel by solving $A\mathbf{x} = \mathbf{0}$:
\[
	\begin{bmatrix}
		2 & -1 & 4 \\
		1 & 3 & -2
	\end{bmatrix} \begin{bmatrix}
		x_1 \\ x_2 \\ x_3
	\end{bmatrix} = \begin{bmatrix}
		0 \\ 0
	\end{bmatrix}
\]

Set up the augmented matrix:
\[
	\left(\begin{array}{ccc|c}
		2 & -1 & 4 & 0 \\
		1 & 3 & -2 & 0
	\end{array}\right)
\]

$R_1 \leftrightarrow R_2$:
\[
	\left(\begin{array}{ccc|c}
		1 & 3 & -2 & 0 \\
		2 & -1 & 4 & 0
	\end{array}\right)
\]

$R_2 \leftarrow R_2 - 2R_1$:
\[
	\left(\begin{array}{ccc|c}
		1 & 3 & -2 & 0 \\
		0 & -7 & 8 & 0
	\end{array}\right)
\]

$R_2 \leftarrow -\frac{1}{7}R_2$:
\[
	\left(\begin{array}{ccc|c}
		1 & 3 & -2 & 0 \\
		0 & 1 & -\frac{8}{7} & 0
	\end{array}\right)
\]

From row 2: $x_2 = \frac{8}{7}x_3$

From row 1: $x_1 = -\frac{10}{7}x_3$

Let $x_3 = t$:
\[
	\begin{bmatrix}
		x_1 \\ x_2 \\ x_3
	\end{bmatrix} = t\begin{bmatrix}
		-\frac{10}{7} \\ \frac{8}{7} \\ 1
	\end{bmatrix} = \frac{t}{7}\begin{bmatrix}
		-10 \\ 8 \\ 7
	\end{bmatrix}
\]

Therefore: $\text{Ker}(T) = \text{span}\left\{\begin{bmatrix} -10 \\ 8 \\ 7 \end{bmatrix}\right\}$
% END QUESTION 5

% BEGIN QUESTION 6
\subsection*{Question 6}
Determine whether the transformation $T : \mathbb{R}^{2} \rightarrow \mathbb{R}^{3}$ defined by:
$$
T(x,y) = (x+y, 2x-y, 3x+2y)
$$
is linear. If yes, find its standard matrix and rank.

\subsection*{Solution}

Check linearity. Let $\mathbf{u} = (x_1, y_1)$ and $\mathbf{v} = (x_2, y_2)$.

$T(\mathbf{u} + \mathbf{v}) = T(x_1 + x_2, y_1 + y_2) = (x_1 + y_1 + x_2 + y_2, 2x_1 - y_1 + 2x_2 - y_2, 3x_1 + 2y_1 + 3x_2 + 2y_2)$

$T(\mathbf{u}) + T(\mathbf{v}) = (x_1 + y_1 + x_2 + y_2, 2x_1 - y_1 + 2x_2 - y_2, 3x_1 + 2y_1 + 3x_2 + 2y_2)$

So $T(\mathbf{u} + \mathbf{v}) = T(\mathbf{u}) + T(\mathbf{v})$

$T(c\mathbf{u}) = T(cx, cy) = (cx + cy, 2cx - cy, 3cx + 2cy) = c(x + y, 2x - y, 3x + 2y) = cT(\mathbf{u})$

Therefore, $T$ is linear.

Apply $T$ to standard basis vectors:
\begin{align*}
	T(1,0) &= (1, 2, 3) \\
	T(0,1) &= (1, -1, 2)
\end{align*}

Standard matrix:
\[
	A = \begin{bmatrix}
		1 & 1 \\
		2 & -1 \\
		3 & 2
	\end{bmatrix}
\]

Row reduce to find rank:
\[
	\begin{bmatrix}
		1 & 1 \\
		2 & -1 \\
		3 & 2
	\end{bmatrix}
\]

$R_2 - 2R_1$:
\[
	\begin{bmatrix}
		1 & 1 \\
		0 & -3 \\
		3 & 2
	\end{bmatrix}
\]

$R_3 - 3R_1$:
\[
	\begin{bmatrix}
		1 & 1 \\
		0 & -3 \\
		0 & -1
	\end{bmatrix}
\]

$-\frac{1}{3}R_2$:
\[
	\begin{bmatrix}
		1 & 1 \\
		0 & 1 \\
		0 & -1
	\end{bmatrix}
\]

$R_3 + R_2$:
\[
	\begin{bmatrix}
		1 & 1 \\
		0 & 1 \\
		0 & 0
	\end{bmatrix}
\]

Two pivot positions, so $\text{rank}(A) = 2$
% END QUESTION 6

\vspace{2em}
\section*{Part C: Cramer's Rule}
\hrulefill
\vspace{1em}

% BEGIN QUESTION 7
\subsection*{Question 7}
Solve the following system using Cramer's Rule:
$$
\begin{cases}
2x + 3y = 7 \\
4x - y = 1
\end{cases}
$$

\subsection*{Solution}

\[
	D = \begin{vmatrix}
		2 & 3 \\
		4 & -1
	\end{vmatrix} = 2(-1) - 3(4) = -14
\]

\[
	D_x = \begin{vmatrix}
		7 & 3 \\
		1 & -1
	\end{vmatrix} = 7(-1) - 3(1) = -10
\]

\[
	D_y = \begin{vmatrix}
		2 & 7 \\
		4 & 1
	\end{vmatrix} = 2(1) - 7(4) = -26
\]

\begin{align*}
	x &= \frac{D_x}{D} = \frac{-10}{-14} = \frac{5}{7} \\
	y &= \frac{D_y}{D} = \frac{-26}{-14} = \frac{13}{7}
\end{align*}

Therefore: $x = \frac{5}{7}, y = \frac{13}{7}$
% END QUESTION 7

% BEGIN QUESTION 8
\subsection*{Question 8}
Use Cramer's Rule to solve the system:
$$
\begin{cases}
x + 2y - z = 4 \\
3x - y + 2z = 1 \\
2x + y + z = 7
\end{cases}
$$

\subsection*{Solution}

\[
	D = \begin{vmatrix}
		1 & 2 & -1 \\
		3 & -1 & 2 \\
		2 & 1 & 1
	\end{vmatrix}
\]
\begin{align*}
	D &= 1\begin{vmatrix}
		-1 & 2 \\
		1 & 1
	\end{vmatrix} - 2\begin{vmatrix}
		3 & 2 \\
		2 & 1
	\end{vmatrix} - 1\begin{vmatrix}
		3 & -1 \\
		2 & 1
	\end{vmatrix} \\
	&= 1(-3) - 2(-1) - 1(5) = -6
\end{align*}

\[
	D_x = \begin{vmatrix}
		4 & 2 & -1 \\
		1 & -1 & 2 \\
		7 & 1 & 1
	\end{vmatrix}
\]

\begin{align*}
	D_x &= 4\begin{vmatrix}
		-1 & 2 \\
		1 & 1
	\end{vmatrix} - 2\begin{vmatrix}
		1 & 2 \\
		7 & 1
	\end{vmatrix} - 1\begin{vmatrix}
		1 & -1 \\
		7 & 1
	\end{vmatrix} \\
	&= 4(-3) - 2(-13) - 1(8) = 6
\end{align*}

\[
	D_y = \begin{vmatrix}
		1 & 4 & -1 \\
		3 & 1 & 2 \\
		2 & 7 & 1
	\end{vmatrix}
\]

\begin{align*}
	D_y &= 1\begin{vmatrix}
		1 & 2 \\
		7 & 1
	\end{vmatrix} - 4\begin{vmatrix}
		3 & 2 \\
		2 & 1
	\end{vmatrix} - 1\begin{vmatrix}
		3 & 1 \\
		2 & 7
	\end{vmatrix} \\
	&= 1(-13) - 4(-1) - 1(19) = -28
\end{align*}

\[
	D_z = \begin{vmatrix}
		1 & 2 & 4 \\
		3 & -1 & 1 \\
		2 & 1 & 7
	\end{vmatrix}
\]

\begin{align*}
	D_z &= 1\begin{vmatrix}
		-1 & 1 \\
		1 & 7
	\end{vmatrix} - 2\begin{vmatrix}
		3 & 1 \\
		2 & 7
	\end{vmatrix} + 4\begin{vmatrix}
		3 & -1 \\
		2 & 1
	\end{vmatrix} \\
	&= 1(-8) - 2(19) + 4(5) = -26
\end{align*}

\begin{align*}
	x &= \frac{D_x}{D} = \frac{6}{-6} = -1 \\
	y &= \frac{D_y}{D} = \frac{-28}{-6} = \frac{14}{3} \\
	z &= \frac{D_z}{D} = \frac{-26}{-6} = \frac{13}{3}
\end{align*}

Therefore: $x = -1, y = \frac{14}{3}, z = \frac{13}{3}$
% END QUESTION 8

% BEGIN QUESTION 9
\subsection*{Question 9}
Determine for what value(s) of k the following system has a unique solution using Cramer's Rule:
$$
\begin{cases}
x + 2y + z = 1 \\
2x + y + z = 2 \\
x + y + kz = 3
\end{cases}
$$

\subsection*{Solution}

\[
	D = \begin{vmatrix}
		1 & 2 & 1 \\
		2 & 1 & 1 \\
		1 & 1 & k
	\end{vmatrix}
\]
\begin{align*}
	D &= 1\begin{vmatrix}
		1 & 1 \\
		1 & k
	\end{vmatrix} - 2\begin{vmatrix}
		2 & 1 \\
		1 & k
	\end{vmatrix} + 1\begin{vmatrix}
		2 & 1 \\
		1 & 1
	\end{vmatrix} \\
	&= k - 1 - 4k + 2 + 1 = -3k + 2
\end{align*}

For unique solution: $D \neq 0$
\[
	-3k + 2 \neq 0 \Rightarrow k \neq \frac{2}{3}
\]

Therefore: The system has a unique solution for all $k \neq \frac{2}{3}$

When $k = \frac{2}{3}$, check by row reduction:

\[
	\left(\begin{array}{ccc|c}
		1 & 2 & 1 & 1 \\
		2 & 1 & 1 & 2 \\
		1 & 1 & \frac{2}{3} & 3
	\end{array}\right)
\]

$R_2 - 2R_1$:
\[
	\left(\begin{array}{ccc|c}
		1 & 2 & 1 & 1 \\
		0 & -3 & -1 & 0 \\
		1 & 1 & \frac{2}{3} & 3
	\end{array}\right)
\]

$R_3 - R_1$:
\[
	\left(\begin{array}{ccc|c}
		1 & 2 & 1 & 1 \\
		0 & -3 & -1 & 0 \\
		0 & -1 & -\frac{1}{3} & 2
	\end{array}\right)
\]

$R_3 - \frac{1}{3}R_2$:
\[
	\left(\begin{array}{ccc|c}
		1 & 2 & 1 & 1 \\
		0 & -3 & -1 & 0 \\
		0 & 0 & 0 & 2
	\end{array}\right)
\]

Last row gives $0 = 2$, contradiction. When $k = \frac{2}{3}$, the system has no solution.
% END QUESTION 9

\end{document}
